\chapter{Transcendental non-Archimedean psh functions}
\index{plurisubharmonic function!Mesquita--Piccione}\index{Monge--Ampere energy!transcendental}
\label{chap:mp-psh}

The tropical spectrum of \Cref{chap:mp-dual-complexes} carries a pluripotential theory once a
K\"ahler class is fixed.  In the projective case positivity can be encoded by
semipositive line-bundle metrics.  A general K\"ahler class need not be the first Chern class of a
line bundle, so positivity is expressed in Bott--Chern cohomology on analytic
test configurations.  This is the approach of \cite{MP24}.

Throughout this chapter, $X$ is a connected compact K\"ahler manifold of dimension
\(n\), \(\alpha\in H^{1,1}_{\mathrm{BC}}(X,\mathbf R)\) is a K\"ahler class,
and
\[
                         V_\alpha:=\int_X\alpha^n.
\tag{60.1}\label{eq:mp-volume}
\]
All measures are Radon measures on the compact Hausdorff space
\(X^{\bethan}\).  We first construct the model theory, then its monotone
completion, and finally explain its exact interface with complex geodesic
rays.  The distinction between these three levels is indispensable.

The order is also conceptual.  Relative nefness is finite-dimensional and
can be tested on a single analytic model.  Energy is obtained by intersecting
classes on common dominating models.  General psh functions then arise by a
monotone completion in which no single model controls the entire net.  The
radial correspondence is a fourth operation: it identifies finite-energy
data with maximal complex rays, but does not turn every psh ray or every DXZ
test curve into model data.

\section{Cohomological models and model potentials}

Let \(\mu:\cX\to X\times\mathbf P^1\) be a smooth snc model and let
\(\pi:\cX\to\mathbf P^1\) be the induced morphism.  Write
\[
          \alpha_{\cX}:=\mu^*\operatorname{pr}_1^*\alpha.
\tag{60.2}\label{eq:mp-pulled-class}
\]
If \(D\in\operatorname{VCar}(\cX)_{\mathbf Q}\) is a rational vertical
Cartier divisor, the pair
\[
                         (\cX,\alpha_{\cX}+[D])
\tag{60.3}\label{eq:mp-cohomological-model}
\]
is the analytic counterpart of a model metric.  Its PL function
\(\phi_D\) is determined by
\[
                     \phi_D(v_E)=b_E^{-1}\ord_E(D)
\tag{60.4}\label{eq:mp-divisor-potential}
\]
for \(\cX_0=\sum_Eb_EE\), and then by affine interpolation on
\(\Delta_{\cX}\).  Pulling \(D\) to a higher model does not change
\(\phi_D\).

\begin{definition}[Model \(\alpha\)-psh functions]
\label{def:mp-model-psh}
A function \(\phi\in\PL(X^{\bethan})\) is \(\alpha\)-psh if, on one
(equivalently every) model on which \(\phi=\phi_D\),
\[
                   \alpha_{\cX}+[D]\in\operatorname{Nef}(\cX/\mathbf P^1).
\tag{60.5}\label{eq:mp-relative-nef}
\]
We write $\PL\cap\PSH(\alpha)$ for this class of functions.
The real model class $\PL_{\mathbf R}\cap\PSH(\alpha)$ is defined
by the same relative nefness condition for real vertical Cartier divisors.
We write $\mathcal H(\alpha)$ for the rational PL functions represented
on some normal model by a relatively K\"ahler class
$\alpha_{\cX}+[D]$.
\end{definition}

If \(X\) is projective and \(\alpha=c_1(L)\), this definition restricts to
the usual algebraic semipositive model metrics.  The operations of addition,
ramified scaling, maxima, and pullback then agree with
\cite[Proposition~3.6]{BJ22}.  Comparing the relatively K\"ahler model
locus with algebraic Fubini--Study functions requires a separate check
of the model category, as explained below.

The parenthetical ``equivalently every'' in
\Cref{def:mp-model-psh} is what makes that definition independent of a
presentation, and its descent half is not formal in the transcendental
K\"ahler setting, so it is recorded separately.

\begin{lemma}[Model independence of relative nefness]
\label{lem:mp-relative-nef-independent}
Let \(\phi\in\PL(X^{\bethan})\).  If \eqref{eq:mp-relative-nef} holds on one
model on which \(\phi=\phi_D\), then it holds on every such model.
\end{lemma}

\begin{proof}
For arbitrary normal models, the model-independence assertion is part of
the imported construction in \cite[Definition~4.1.2, arXiv version~2]{MP24}.
We give a direct proof on the smooth snc models used in the intersection
formulas below.  It is a consequence of
\cite[Theorem~4.3(iii), Proposition~3.3(ii)]{DP04}.

Any two such models admit a common smooth snc domination.  On the common
model, the two pulled-back vertical divisors have the same coefficient
$b_F\phi(v_F)$ at every central prime $F$, so the represented classes
agree.  It suffices to compare a smooth snc model $\cX$ with a smooth
snc modification $q:\mathcal Y\to\cX$, writing
$\beta=\alpha_{\cX}+[D]$.  Pullback preserves nefness on each fibre,
which proves one implication.

For the converse, assume $q^*\beta$ is relatively nef.  Fix a central
component $E$ of $\cX$, an irreducible analytic subset $Z\subseteq E$
of dimension $d\geq1$, and a K\"ahler form $\gamma$ on $E$.
Choose an irreducible component $W$ of $q^{-1}(Z)$ which maps onto $Z$.
Such a component exists because $q$ is proper and surjective.  An embedded
resolution of $W$ in $\mathcal Y$ gives a compact K\"ahler manifold
$\widetilde W$.  Put
$r=\dim\widetilde W\geq d$, and let
$g:\widetilde W\to E$ be the induced map.  Relative nefness upstairs
implies that $g^*(\beta|_E)$ is nef: the resolution map factors through
a central component of $\mathcal Y$.

Choose a K\"ahler form $\kappa$ on $\widetilde W$.  The closed positive
current $g_*(\kappa^{r-d})$ has bidimension $(d,d)$ and support $Z$.
The support theorem therefore gives
\[
                 g_*(\kappa^{r-d})=c[Z],\qquad c>0.
\]
The strict positivity of $c$ follows by pairing with $\gamma^d$:
$g$ has rank $d$ on a nonempty open subset, so
$\int_{\widetilde W}\kappa^{r-d}\wedge g^*\gamma^d>0$.
Projection formula now yields
\[
 c\int_Z\beta|_E\wedge\gamma^{d-1}
 =\int_{\widetilde W}
      g^*(\beta|_E)\wedge g^*\gamma^{d-1}\wedge\kappa^{r-d}
 \geq0.
\]
The last inequality follows by approximating the nef class
$g^*(\beta|_E)$ by K\"ahler classes; $g^*\gamma$ is semipositive.
The numerical criterion gives nefness of $\beta|_E$.  The finite-union
assertion gives nefness on the reduced central fibre; the other fibres
carry the original K\"ahler class $\alpha$.  This proves relative nefness
downstairs.  Zero-dimensional fibres require no numerical test.

Choosing a component of the full inverse image is essential here: a
subvariety contained in the exceptional image need not have a strict
transform.  The proof uses every analytic dimension and every K\"ahler
test class on $E$.
\end{proof}

The corresponding statement for $\mathcal H(\alpha)$ is false:
a relatively K\"ahler class can cease to be relatively K\"ahler after
pullback to a higher model, although it remains relatively nef.  The clause defining \(\mathcal H(\alpha)\) in
\Cref{def:mp-model-psh} is therefore to be read as ``represented by a
relatively K\"ahler class on \emph{some} model'', and only the nef cone
carries the ``one (equivalently every)'' quantifier.

\begin{warning}{Model categories in Fubini--Study comparisons}
\label{warn:mp-h-not-fs}
Suppose that $X$ is smooth projective, $L$ is an ample
$\mathbf Q$-line bundle, and $\alpha=c_1(L)$.  On a smooth projective
algebraic model $\mu:\cX\to X\times\mathbf P^1$, a rational vertical
divisor $D$ does give a line-bundle class.  Choose a positive integer $m$ so that
$L^{\otimes m}$ is a line bundle and $mD$ is integral Cartier.  Then
\[
 \mathcal M=
 \mu^*\operatorname{pr}_1^*L^{\otimes m}\otimes\mathcal O_{\cX}(mD),
 \qquad
 m^{-1}c_1(\mathcal M)=\alpha_{\cX}+[D].
\]
If the latter class is relatively K\"ahler, its restriction to every
positive-dimensional subvariety of a fibre has positive top
self-intersection.  Fibrewise Nakai--Moishezon makes $\mathcal M$ ample on every fibre.
The fibre criterion \cite[Tag~0D2S]{stacks-project} then makes it
relatively ample.
Consequently $\phi_D$ is an algebraic Fubini--Study function; see
\cite[Corollary~2.32]{BJ22}.

The model category still matters.  In
\cite[Definition~3.8 and Example~3.9]{BJ22}, the locus
$\mathcal H^{\mathrm{dom}}(L)\subseteq\mathcal H(L)$ consists of those
Fubini--Study functions whose ample, integrally closed test configuration
dominates the trivial one.  This is the algebraic locus directly described
by relative ampleness of $\alpha_{\cX}+[D]$ on a dominating model.
Passing to a common resolution preserves nefness, but can destroy
relative ampleness.  Comparisons of the strictly positive loci must
therefore specify the background class and the allowed models.

The full psh comparison in \Cref{thm:pseff-projective-full-mp-bj} uses
relative nefness on common models and decreasing limits, without
identifying strictly positive model loci.
\end{warning}

\begin{lemma}
\label{lem:mp-relative-kahler-extension}
Let $\pi:\cX\to\mathbf P^1$ be a smooth snc model and let
$\omega_B$ be a K\"ahler form on $\mathbf P^1$.  If a real
Bott--Chern class $\xi$ on $\cX$ is K\"ahler on every reduced fibre,
then $\xi+C\pi^*[\omega_B]$ is K\"ahler on $\cX$ for some $C\geq0$.
\end{lemma}

\begin{proof}
Put $\eta=\pi^*\omega_B$ and choose a smooth representative
$\theta$ of $\xi$.  By \cite[Proposition~3.3(i)--(ii)]{DP04}, near
each fibre there is a representative $\theta+\ddc\psi_j$ which is positive on an
ambient neighbourhood; for the central fibre apply the union assertion
to its finitely many components first.  Properness provides a finite
cover of the base whose inverse images lie in these neighbourhoods.
Choose a smooth partition of unity $\chi_j$ on the base subordinate to
that cover.  The form
\[
 \theta'=\theta+\ddc\sum_j(\chi_j\circ\pi)\psi_j
\]
is positive on $\ker\mathrm{d}\pi$: the terms involving
$\mathrm{d}(\chi_j\circ\pi)$ or $\ddc(\chi_j\circ\pi)$ vanish there,
so its restriction is the weighted sum of the positive local representatives.  The form $\eta$ is semipositive
and vanishes exactly on that kernel.  Fixing a Hermitian metric on
$\cX$, compactness of its unit tangent bundle gives $\theta'+C\eta>0$ for some $C$.  Indeed, if no
such $C$ existed, unit vectors violating positivity for $C\to\infty$
would have a subsequence tending to a vector in the kernel on which
$\theta'$ is nonpositive.  This contradicts the established positivity.
\end{proof}

The positive-rational operation below is the scaling action
\[
             (q\cdot\phi)(v):=q\,\phi(q^{-1}v),
\tag{60.5a}\label{eq:mp-scaling-action}
\]
not pointwise scalar multiplication.  It is induced by ramified base change
and fixes \(v_{\triv}\).

\begin{proposition}\label{prop:mp-model-psh-stability}
The following operations preserve the indicated model psh classes:
\begin{enumerate}
\item \(\phi\mapsto\phi+c\) preserves
      \(\PL\cap\PSH(\alpha)\) for \(c\in\mathbf Q\); for
      \(c\in\mathbf R\), it preserves the real model class
      \(\PL_{\mathbf R}\cap\PSH(\alpha)\);
\item if \(q\in\mathbf Q_{>0}\), then
      \(q\cdot\phi\in\PL\cap\PSH(\alpha)\);
\item \(\max\{\phi,\psi\}\in\PL\cap\PSH(\alpha)\) whenever
      \(\phi,\psi\in\PL\cap\PSH(\alpha)\);
\item if \(f:Y\to X\) is finite holomorphic, then
      \(f^*\phi\in\PL\cap\PSH(f^*\alpha)\).
\end{enumerate}
\par
\end{proposition}

\begin{proof}
Adding \(c\) replaces \(D\) by \(D+c\cX_0\), which has the same relative
class.  The divisor remains rational when \(c\in\mathbf Q\), and the same
argument applies to real vertical divisors when \(c\in\mathbf R\).  For
integral \(q\), normalize the base change
\(t\mapsto t^q\); pullback of the relatively nef class represents
\(q\cdot\phi\).  If \(q=a/b\in\mathbf Q_{>0}\), apply the integral case to
\(a\), and descend the assertion through the finite base change of degree
\(b\), using that a class is nef if and only if its pullback by a finite
surjective map is nef.  This is why the background class remains
\(\alpha\).  Finite
pullback preserves nefness.

For the maximum, choose a common smooth snc model $\cX$ on which
$\phi=\phi_D$ and $\psi=\phi_E$.  Choose an integer $c\geq0$ and
a positive integer $m$ such that
\[
 F=m(c\cX_0-D),\qquad G=m(c\cX_0-E)
\]
are effective integral Cartier divisors.  Let
\[
 \mathfrak a=\mathcal O_{\cX}(-F)+\mathcal O_{\cX}(-G),
 \qquad
 \mathfrak a\mathcal O_{\mathcal Y}=\mathcal O_{\mathcal Y}(-H)
\]
on an equivariant log resolution $q:\mathcal Y\to\cX$.  We may take
$\mathcal Y$ snc.  The valuation identity for a sum of ideals gives
\[
 \max\{\phi,\psi\}=\phi_{D_{\max}},
 \qquad D_{\max}=c\mathcal Y_0-m^{-1}H.
\]
In particular, the minus signs in the two ideal sheaves are essential.

Fix K\"ahler forms $\Omega$ on $\cX$ and $\omega_B$ on
$\mathbf P^1$, and write $\eta=\pi^*\omega_B$ for the pulled-back
class.  For each $\varepsilon>0$, the classes
$m(\alpha_{\cX}+[D])+\varepsilon\Omega$ and
$m(\alpha_{\cX}+[E])+\varepsilon\Omega$ are K\"ahler on every
reduced fibre.  Here nefness on the smooth central components and
\cite[Proposition~3.3(ii)]{DP04} give positivity on their union.
By \Cref{lem:mp-relative-kahler-extension}, a common $C_\varepsilon\geq0$
makes both classes K\"ahler after adding $C_\varepsilon\eta$.  Thus
\[
 \beta_\varepsilon
 =m(\alpha_{\cX}+c[\cX_0])+\varepsilon\Omega
       +C_\varepsilon\eta
\]
has the property that $\beta_\varepsilon-[F]$ and
$\beta_\varepsilon-[G]$ admit smooth positive representatives.
The class $\beta_\varepsilon$ is itself K\"ahler, since
$\alpha_{\cX}$ and $[\cX_0]$ have semipositive representatives.

Choose a K\"ahler representative $\sigma_\varepsilon$ of
$\beta_\varepsilon$ and metrics on $\mathcal O(F)$ and
$\mathcal O(G)$ whose curvature forms $\theta_F,\theta_G$ satisfy
$\sigma_\varepsilon-\theta_F>0$ and
$\sigma_\varepsilon-\theta_G>0$.  For the canonical sections $s_F,s_G$,
the convention \eqref{eq:complex-ddc-normalization} gives
$\ddc\log|s_F|^2=[F]-\theta_F$ and the corresponding identity for $G$.
Consequently
\[
 u=\max\{\log|s_F|^2,\log|s_G|^2\}
\]
is $\sigma_\varepsilon$-psh and has singularity type
$\log|\mathfrak a|^2$.  After pullback to $\mathcal Y$, subtracting
$\log|f_H|^2$ for a local equation $f_H$ of $H$ leaves a locally
bounded potential: the two remaining ideal generators have no common zero.
The bounded-potential regularization argument in
\cite[Proposition~5.0.1, arXiv version~2]{MP24} therefore gives
\[
 q^*\beta_\varepsilon-[H]
 =m(\alpha_{\mathcal Y}+[D_{\max}])
      +\varepsilon q^*\Omega+C_\varepsilon q^*\eta
 \in\operatorname{Nef}(\mathcal Y).
\]
Restrict this class to each smooth central component of $\mathcal Y$.
The base term vanishes there.  Letting $\varepsilon\downarrow0$ and
using closedness of the nef cone proves nefness of
$\alpha_{\mathcal Y}+[D_{\max}]$ on that component.  The other fibres
carry $\alpha$, so this class is relatively nef.  The resolution and
$D_{\max}$ were fixed throughout; no uniform bound on
$C_\varepsilon$ is needed.  This proves the maximum assertion directly
for relatively nef model classes.
\end{proof}

\begin{remark}[Why rationality matters]
The fixed-class scaling above is precisely
\cite[Proposition~4.1.4(2)]{MP24}; it is not the elementary operation
\(\alpha_{\cX}+[D]\mapsto\alpha_{\cX}+[qD]\) on a fixed model.  It uses the
change of Gauss normalization under a ramified base change.  The statement
is made only for \(q\in\mathbf Q_{>0}\); we do not infer arbitrary real
scaling at the PL level.
\end{remark}

\begin{theorem}\label{thm:mp-model-closedness}
Let \((\phi_\lambda)\) be a net in
\(\PL\cap\PSH(\alpha)\), and suppose that
\(\phi_\lambda(v)\to\phi(v)\) for every \(v\in X^{\divis}\), where
\(\phi\in\PL(X^{\bethan})\).  Then
\(\phi\in\PL\cap\PSH(\alpha)\).
\end{theorem}

\begin{proof}[Imported analytic-cycle core]
The imported core is \cite[Theorem~4.1.5, arXiv version~2]{MP24}, whose
statement is exactly the present one: $X$ is compact K\"ahler, $\alpha$
is K\"ahler, $\phi$ is PL, and a net of model $\alpha$-psh functions
converges to $\phi$ at every rational divisorial valuation.  We give the
relative numerical reduction in detail and then recall the analytic-cycle
argument.

Choose a smooth snc model $\cX$ representing $\phi$.  For each
$\lambda$, pass to a common domination of $\cX$ and one model
representing $\phi_\lambda$; no single model dominating the entire net is
assumed.  Write $\beta=\alpha_{\cX}+[D]$.  The numerical reduction in the
imported theorem tests
\[
 [Y]\cdot\beta\cdot\gamma^{d-1}\geq0
\tag{60.6}\label{eq:mp-nef-test}
\]
for every positive-dimensional irreducible $Y$ in a central component,
$d=\dim Y$, and every K\"ahler class $\gamma$ on $\cX$.  We first justify
this relative reduction, including the fact that the test classes live on
the total space.

\emph{The relative numerical reduction.}
The underlying absolute criterion is
\cite[Theorem~4.3(iii), Corollary~0.4]{DP04}.  Fix K\"ahler forms
$\Omega$ on $\cX$ and $\omega_B$ on $\mathbf P^1$, and put
$\eta=\pi^*\omega_B$; the same symbols denote their cohomology classes.
If $n=0$, relative nefness is automatic.  Assume $n\geq1$.

Suppose now that \eqref{eq:mp-nef-test} holds and that $a>0$ is such that
$\beta+a\Omega$ is K\"ahler on every fibre.  Such an $a$ exists initially
because $\Omega$ is K\"ahler on the compact total space.
\Cref{lem:mp-relative-kahler-extension} gives $C_a\geq0$ for which
$\beta+a\Omega+C_a\eta$ is K\"ahler.  Consequently every class
\[
 \gamma_s=s\beta+a\Omega+C_a\eta,\qquad 0\leq s\leq1,
\]
is K\"ahler, being a convex combination of that class and
$a\Omega+C_a\eta$.

For each $1\leq d\leq n$, use the polynomial identity
\cite[Equation~(4.4) and its proof, pp.~1267--1268]{DP04}, in the
normalization
\[
 (y-\delta x)^d-(1-\delta)^d x^d
 =d(y-x)\int_0^1 B_d(s,\delta)((1-s)x+sy)^{d-1}\,\mathrm{d}s,
\]
where $B_d$ is a polynomial of degree at most $d-1$ in $s$ and
$B_d(s,0)=1$.  To check this normalization and existence, divide the
left-hand side by $d(y-x)$ and compare its coefficients in the basis
$x^{d-1-k}y^k$.  The resulting moments determine a unique polynomial
in $s$, because integration gives a nondegenerate pairing on
polynomials of degree at most $d-1$.  Its coefficients depend
polynomially on $\delta$.  At $\delta=0$, the fundamental theorem of
calculus gives $B_d(s,0)=1$ by uniqueness.  Choose one
$0<\delta_0<1$ such that all these polynomials are positive for
$0\leq s\leq1$ and $0\leq\delta\leq\delta_0$.
Restricting the identity to a central $d$-dimensional subvariety $Y$
eliminates $\eta$ and gives
\begin{align*}
 \int_Y\bigl(\beta+(1-\delta)a\Omega\bigr)^d
 &-\bigl((1-\delta)a\bigr)^d\int_Y\Omega^d\\
 &=d\int_0^1 B_d(s,\delta)
       \left(\int_Y\beta\wedge\gamma_s^{d-1}\right)
          \,\mathrm{d}s\geq0.
\end{align*}
In particular the top intersection on the first line is strictly
positive on every such $Y$.  On each smooth central component the
segment $\beta+(1-\delta)a\Omega$, $0\leq\delta\leq\delta_0$, starts
in the K\"ahler cone and stays in the cone of classes positive on all
analytic cycles.  By \cite[Theorem~0.1]{DP04} its whole segment is
K\"ahler on that component.  Away from the central fibre, $\beta$
restricts to the original K\"ahler class $\alpha$, so the same
conclusion holds on every other fibre.  The union assertion above handles
the reduced central fibre.  Thus $\beta+(1-\delta_0)a\Omega$ is again
K\"ahler on every fibre.

Iterating with $a_m=(1-\delta_0)^m a_0$ and letting $m\to\infty$ proves
that $\beta$ is nef on every fibre, hence relatively nef.  The number
$\delta_0$ is fixed throughout; the constants $C_{a_m}$ may vary with
$m$, since the hypothesis tests every ambient K\"ahler class.
Conversely, relative nefness implies \eqref{eq:mp-nef-test} by restricting
to the central components and taking limits of positive perturbations.

\emph{The analytic-cycle limit.}
It remains to obtain \eqref{eq:mp-nef-test} from the divisorial convergence.
The cited analytic-cycle argument first replaces $Y$ by an
effective $\mathbf C^*$-invariant limiting cycle and treats its components
separately.  For an invariant component, a fixed blowup carries a positive
exceptional current whose pushforward is a strictly positive multiple of
$[Y]$.  On a common domination with the model of $\phi_\lambda$, relative
nefness makes the corresponding intersection nonnegative.  Pushing back
to the fixed blowup expresses it as an affine function of finitely many
values $\phi_\lambda(v_E)$; the constant term comes from the fixed
background class.  Divisorial convergence therefore gives
\eqref{eq:mp-nef-test}.  The cycle limit, the positivity of the multiple,
and the projection-formula calculation remain the imported analytic core.
The common domination is chosen separately for each index $\lambda$.
\end{proof}

\section{Model Monge--Amp\`ere measures}

Let \(\beta\in H^{1,1}_{\mathrm{BC}}(X,\mathbf R)\) satisfy
\(V_\beta:=\int_X\beta^n>0\), and let \(\phi_D\) be represented on an snc
model \(\cX\).  Define the signed atomic measure
\[
 \MA_\beta(\phi_D)
   :=\sum_{E\subset\cX_0}
       \frac{b_E}{V_\beta}
       \left((\beta_{\cX}+[D])|_E\right)^n\delta_{v_E}.
\tag{60.7}\label{eq:mp-model-ma}
\]
The same formula defines $\MA_\beta(\phi)$ for
$\phi\in\PL_{\mathbf R}$, using a real vertical divisor on a common
model.  The sum is over irreducible components.  To check independence of the
model, let $q:\mathcal Y\to\cX$ be a smooth snc domination.  A strict
transform of a central component has the same multiplicity and divisorial
valuation, and its top intersection is unchanged by projection formula.
For every other central component $F$ of $\mathcal Y$, the image
$q(F)$ has dimension less than $n$, so
\[
 \int_F\bigl(q^*(\beta_{\cX}+[D])\bigr)^n=0.
\]
Thus the extra components contribute zero to the measure.  This is the
construction in \cite[Section~4.2.1, arXiv version~2]{MP24}.

\begin{proposition}\label{prop:mp-model-ma-mass}
For $V_\beta>0$ and $\phi\in\PL_{\mathbf R}$, the signed measure
$\MA_\beta(\phi)$ has total mass one.  If $\beta$ is K\"ahler and
$\phi\in\PL_{\mathbf R}\cap\PSH(\beta)$, then it is a probability measure.
\end{proposition}

\begin{proof}
The total mass is
\[
 \frac1{V_\beta}\sum_Eb_E
     (\beta_{\cX}+[D])^n\cdot E
 =\frac1{V_\beta}(\beta_{\cX}+[D])^n\cdot\cX_0=1.
\]
Indeed, $\cX_0$ is cohomologous to a general fibre, where $D$ restricts
trivially and the remaining top intersection is $V_\beta$.
If $\phi$ is model $\beta$-psh, the class
$(\beta_{\cX}+[D])|_E$ is nef on each smooth central component.
Its top self-intersection is nonnegative, by approximation with K\"ahler
classes.  Since $b_E>0$ and $V_\beta>0$, every atomic coefficient is
nonnegative.  In dimension zero the top intersection is one.
\end{proof}

\begin{example}[A deformation to the normal cone]
\label{ex:mp-dnc}
Assume $n\geq1$.  Let $Z\subsetneq X$ be a
nonempty connected closed complex submanifold and put
$\cX=\operatorname{Bl}_{Z\times\{0\}}(X\times\mathbf P^1)$.
The central fibre has exactly two irreducible components: the strict
transform $E_0$ and the exceptional component $E_1$.
For $D=-cE_1$ with $c\in\mathbf Q_{>0}$ sufficiently small,
\(\alpha_{\cX}+[D]\) is relatively K\"ahler.  Thus
\(\phi_D\) is affine on the interval \(\Delta_{\cX}\), and
\[
 \begin{gathered}
 \MA_\alpha(\phi_D)
   =m_0(c)\delta_{v_{E_0}}+m_1(c)\delta_{v_{E_1}},\\
 m_i(c)\geq0,\qquad m_0(c)+m_1(c)=1.
 \end{gathered}
\tag{60.8}\label{eq:mp-dnc-ma}
\]
This construction is meaningful for a transcendental \(\alpha\); no
polarization of \(X\) is required.

Every coefficient is computable once \(Z\) is a point.  Let
\(Z=\{p\}\) with \(X\) of dimension \(n\), so that
\(E_0=\operatorname{Bl}_pX\), \(E_1\cong\mathbf P^n\), and
\(b_0=b_1=1\).  Write \(A=\alpha_{\cX}\) and \(B=A-c[E_1]\).  Since
\(E_1\) maps to a point, \(A|_{E_1}=0\), while
\(E_1|_{E_1}=-h\) for the hyperplane class \(h\) on \(E_1\); hence
\(B|_{E_1}=ch\) and \((B|_{E_1})^n=c^n\).  For $n\geq2$, the restriction to the strict transform is
$B|_{E_0}=\rho^*\alpha-c\,e$, where
$\rho:\operatorname{Bl}_pX\to X$ and $e$ is its exceptional divisor.
Since $\rho^*\alpha\cdot e=0$ and $\int_{E_0}e^n=(-1)^{n-1}$,
the top intersection is $V_\alpha-c^n$.  When $n=1$, the blowup
$\rho$ is the identity and $E_0\cap E_1$ is the point $p$; thus
$\deg(B|_{E_0})=V_\alpha-c$, giving the same formula.  Therefore
\[
 \MA_\alpha(\phi_D)
 =\Bigl(1-\frac{c^n}{V_\alpha}\Bigr)\delta_{v_{E_0}}
  +\frac{c^n}{V_\alpha}\,\delta_{v_{E_1}},
\]
which is a probability measure, as
\Cref{prop:mp-model-ma-mass} requires.

Taking \(X=\mathbf P^2\) and \(\alpha=c_1(\cO(1))\), so that
\(n=2\) and \(V_\alpha=1\), the remaining invariants are numbers.  All
mixed terms in \(B^{3}\) vanish --- \(A^3=0\) because \(A\) is pulled back
from a surface, and \(A\cdot E_1=0\) on \(E_1\) --- leaving
\(B^3=-c^3E_1^3=-c^3\).  Hence
\[
 \MA_\alpha(\phi_D)=(1-c^2)\delta_{v_{E_0}}+c^2\delta_{v_{E_1}},
 \qquad
 E_\alpha(\phi_D)=-\frac{c^3}{3},
 \qquad
 J_\alpha(\phi_D)=\frac{c^3}{3},
\]
the last because \(B\cdot A^2=0\).  Finally
\(K_{\cX/(X\times\mathbf P^1)}=nE_1\) gives
\(A^{\mathrm{tot}}(v_{E_1})=n+1\) and
\(A^{\mathrm{tot}}(v_{E_0})=1\), so the reduced discrepancies are
\(A^{\NA}_X(v_{E_1})=2\) and \(A^{\NA}_X(v_{E_0})=0\), and the entropy is
\[
 \Ent^{\NA}_\alpha(\phi_D)
 =\int A^{\NA}_X\,\rd\MA_\alpha(\phi_D)=2c^2 .
\]
This one configuration instantiates
\eqref{eq:mp-divisor-potential}, \eqref{eq:mp-model-ma},
\eqref{eq:mp-energy-j} and the reduced-versus-total shift
\eqref{eq:mp-total-reduced-mabuchi} in closed form.
\end{example}

\section{Mixed energy and the Zariski inequality}

For \(\beta_i\in H^{1,1}_{\mathrm{BC}}(X,\mathbf R)\) and model functions
\(\phi_i=\phi_{D_i}\) represented on one model, define
\[
 (\beta_0,\phi_0)\cdots(\beta_n,\phi_n)
 :=(\beta_{0,\cX}+[D_0])\cdots(\beta_{n,\cX}+[D_n]).
\tag{60.9}\label{eq:mp-mixed-pairing}
\]
The right side is the Bott--Chern intersection number on the compact
\((n+1)\)-fold \(\cX\).  It is symmetric, multilinear, and independent of
the common model.

\begin{proposition}\label{prop:mp-ma-testing}
If \(\psi\in\PL_{\mathbf R}\), then
\[
 \frac1{V_\alpha}(0,\psi)(\alpha,\phi)^n
     =\int_{X^{\bethan}}\psi\,\MA_\alpha(\phi).
\tag{60.10}\label{eq:mp-ma-testing}
\]
More generally,
\[
 (0,\psi)(\beta_1,\phi_1)\cdots(\beta_n,\phi_n)
 =\sum_{E\subset\cX_0}b_E\psi(v_E)
      \prod_{j=1}^n(\beta_{j,\cX}+[D_j])|_E .
\tag{60.11}\label{eq:mp-mixed-testing}
\]
\end{proposition}

\begin{proof}
Write the divisor representing \(\psi\) as
\(\sum_Eb_E\psi(v_E)E\) after passing to a common model.  Expanding its
intersection with the remaining \(n\) classes yields
\eqref{eq:mp-mixed-testing}; the special case is precisely
\eqref{eq:mp-model-ma}.
\end{proof}

For $n\geq1$, the vertical intersection form satisfies the Zariski inequality
\[
             (0,\psi)^2(\beta_2,\phi_2)\cdots(\beta_n,\phi_n)\leq0
\tag{60.12}\label{eq:mp-zariski}
\]
for $\psi\in\PL_{\mathbf R}$ when the remaining pairs are relatively
nef.  The proof is a finite-dimensional intersection calculation;
see \cite[Lemma~4.2.6, arXiv version~2]{MP24}.
Choose a common smooth snc model with $\psi=\phi_D$ and
$\phi_k=\phi_{D_k}$, and write
\[
 \cX_0=\sum_i b_iE_i,\qquad
 D=\sum_i b_i a_iE_i,\qquad a_i=\psi(v_{E_i}).
\]
Put
\[
 \Gamma=\prod_{k=2}^n(\beta_{k,\cX}+[D_k]),\qquad
 Q_{ij}=E_i\cdot E_j\cdot\Gamma,
\]
where the product is $1$ if $n=1$.  For $i\ne j$, the intersection
$E_i\cap E_j$ is smooth, possibly empty or disconnected, and the
classes in $\Gamma$ restrict to nef classes there.  Approximating those classes
by K\"ahler classes gives $Q_{ij}\geq0$.  Moreover,
\[
 \sum_j b_jQ_{ij}=E_i\cdot\cX_0\cdot\Gamma=0,
\]
since $\mathcal O(\cX_0)$ restricts trivially to $E_i$.
Expanding the square and using these weighted row sums gives
\[
 D^2\cdot\Gamma
 =-\sum_{i<j}b_i b_jQ_{ij}(a_i-a_j)^2\leq0.
\]
This proves \eqref{eq:mp-zariski}, including arbitrary real coefficients
and vanishing intersection weights.

\begin{definition}[Energy and $J$]
\label{def:mp-energy-j}
For $\phi\in\PL_{\mathbf R}$, set
\[
 \begin{aligned}
 E_\alpha(\phi)
   &:=\frac1{(n+1)V_\alpha}(\alpha,\phi)^{n+1},\\
 J_\alpha(\phi)
   &:=V_\alpha^{-1}(\alpha,\phi)(\alpha,0)^n-E_\alpha(\phi).
 \end{aligned}
\tag{60.13}\label{eq:mp-energy-j}
\]
\end{definition}

These expressions are polynomials on the real space of vertical divisors
of any fixed model.  In particular, they are defined along real affine
segments, including parameters for which the divisor is not rational.
On rational PL functions, the energy agrees with
\cite[Definition~4.2.9, arXiv version~2]{MP24}.

\begin{proposition}\label{prop:mp-energy-first-variation}
For $\phi,\psi\in\PL_{\mathbf R}$,
\[
 \left.\frac{\rd}{\rd t}\right|_{t=0}
 E_\alpha(\phi+t\psi)
       =\int\psi\,\MA_\alpha(\phi).
\tag{60.14}\label{eq:mp-energy-first-variation}
\]
The energy is increasing and concave on the real model $\alpha$-psh
class.  For arbitrary $\phi,\psi\in\PL_{\mathbf R}$, one also has
\[
 E_\alpha(\phi)-E_\alpha(\psi)
 =\frac1{(n+1)V_\alpha}
   \sum_{j=0}^n
   (0,\phi-\psi)(\alpha,\phi)^j(\alpha,\psi)^{n-j}.
\tag{60.15}\label{eq:mp-energy-cocycle}
\]
\end{proposition}

\begin{proof}
Put $B_t=(\alpha,\phi+t\psi)$.  Multilinearity gives
\[
 \frac{\rd}{\rd t}B_t^{n+1}
      =(n+1)(0,\psi)B_t^n.
\]
After division by $(n+1)V_\alpha$,
\eqref{eq:mp-ma-testing} identifies the right-hand side at $t=0$ with
$\int\psi\,\MA_\alpha(\phi)$, proving
\eqref{eq:mp-energy-first-variation}.  The direction $\psi$ need not
be psh for this polynomial identity.

Now let $\phi,\psi$ be real model $\alpha$-psh functions and put
$u_t=(1-t)\phi+t\psi$ for $0\leq t\leq1$.
Convexity of the relative nef cone makes $u_t$ model $\alpha$-psh.
For $n\geq1$, a second differentiation gives
\[
 \frac{\rd^2}{\rd t^2}E_\alpha(u_t)
   =\frac{n}{V_\alpha}(0,\psi-\phi)^2(\alpha,u_t)^{n-1}\leq0
\]
by \eqref{eq:mp-zariski}.  If $n=0$, the defining polynomial is affine.
This proves concavity in every dimension.  If $\phi\leq\psi$, then
\[
 \frac{\rd}{\rd t}E_\alpha(u_t)
   =\int(\psi-\phi)\,\MA_\alpha(u_t)\geq0
\]
by \Cref{prop:mp-model-ma-mass}, which proves monotonicity.  This
uses positivity of the measure and of the difference as a function;
no psh property of the difference is needed.

Finally, multilinearity gives
\[
 (\alpha,\phi)^{n+1}-(\alpha,\psi)^{n+1}
 =\sum_{j=0}^n
   (0,\phi-\psi)(\alpha,\phi)^j(\alpha,\psi)^{n-j}.
\]
Dividing by $(n+1)V_\alpha$ proves
\eqref{eq:mp-energy-cocycle}.
\end{proof}

\section{The monotone completion}

\begin{definition}[Transcendental non-Archimedean psh functions]
\label{def:mp-general-psh}
An upper semicontinuous function
\(\phi:X^{\bethan}\to\mathbf R\cup\{-\infty\}\) is
\(\alpha\)-psh if it is not identically \(-\infty\) and there exists a
decreasing net
\[
        \phi_\lambda\in\PL\cap\PSH(\alpha),
        \qquad \phi_\lambda\downarrow\phi
\tag{60.16}\label{eq:mp-psh-decreasing-net}
\]
pointwise on all of \(X^{\bethan}\).  The resulting space is
\(\PSH(\alpha)\).
\end{definition}

Nets are part of the definition: the model category and the topology of the
tropical spectrum need not be countably generated.  The next result gives
a sequence when the target function has a model representative; it does
not assert countable cofinality of the whole model category.

\begin{proposition}
\label{prop:mp-real-model-approximation}
Let $\phi=\phi_D$ be a real model $\alpha$-psh function.
There is a sequence $\psi_j\in\mathcal H(\alpha)$ decreasing uniformly
to $\phi$.  The divisors representing $\psi_j$ may be chosen on one
smooth snc model dominating the model of $D$, with their coefficients
converging to those of the pullback of $D$.  In particular, every real
model $\alpha$-psh function belongs to $\PSH(\alpha)$ as defined above.
\end{proposition}

\begin{proof}
Pass to a smooth snc model $\mu:\mathcal Y\to X\times\mathbf P^1$
obtained by equivariant blowups and dominating the model of $D$.
The cofinality and relative-ampleness constructions used here are those
of \cite[Remark~2.5.2 and the proof of Theorem~2.6.4,
arXiv version~2]{MP24}.  They give an integral vertical divisor $F$
such that $\mathcal O_{\mathcal Y}(F)$ is $\mu$-ample.  Concretely, on
a sequence of blowups one combines the negatives of the exceptional
divisors, choosing each earlier coefficient sufficiently large compared
with the later ones.  We retain the notation $D$ for its pullback.

There is a rational vertical divisor $A$ on $\mathcal Y$ such that
$\alpha_{\mathcal Y}+[A]$ is relatively K\"ahler.  To see this, choose
a smooth curvature form $\theta_F$ of $F$ which is positive on
$\ker\mathrm d\mu$.  Such a form is obtained from local relative
projective embeddings; metrics can be patched by a partition of unity
on the base without changing their positivity on that kernel.
For K\"ahler forms $\omega\in\alpha$ and $\omega_B$ on $\mathbf P^1$,
compactness of the unit tangent bundle shows that
\[
 \mu^*(\operatorname{pr}_1^*\omega+
             \operatorname{pr}_2^*\omega_B)+s\theta_F>0
\]
for all sufficiently small $s>0$.  Indeed, vectors violating this
inequality as $s\downarrow0$ would converge to a unit vector in
$\ker\mathrm d\mu$, where $\theta_F$ is positive.  Choose such an
$s\in\mathbf Q$ and put $A=sF$.  Restricting to the fibres over
$\mathbf P^1$ proves the assertion.

For $0<t\leq1$, the divisor $D_t=(1-t)D+tA$ represents a relatively
K\"ahler class: on every central component it is the sum of a nef
class and a positive multiple of a K\"ahler class.  The same holds on
their union by \cite[Proposition~3.3(ii)]{DP04}.  Relative K\"ahlerness
is open in the finite-dimensional space of vertical coefficients on
$\mathcal Y$.  This follows from openness on its finitely many smooth
central components; every other fibre has the fixed class $\alpha$.
Consequently there are rational vertical divisors $B_j$ converging
coefficientwise to $D$ such that $\alpha_{\mathcal Y}+[B_j]$ is
relatively K\"ahler and
\[
        \|\phi_{B_j}-\phi\|_\infty<\varepsilon_j,
        \qquad \varepsilon_j=4^{-j}.
\]
Here coefficient convergence implies uniform convergence, since the
model functions factor through the fixed finite dual complex and are
affine on each face.

Set $\psi_j=\phi_{B_j}+2\varepsilon_j$.  Rational translation preserves
$\mathcal H(\alpha)$, and
\[
 \phi+\varepsilon_j<\psi_j<\phi+3\varepsilon_j,
 \qquad 3\varepsilon_{j+1}<\varepsilon_j.
\]
Thus $\psi_{j+1}<\psi_j$ everywhere and $\psi_j\to\phi$ uniformly.
Their divisors are $B_j+2\varepsilon_j\mathcal Y_0$, whose coefficients
still converge to those of $D$.
\end{proof}

\begin{proposition}
\label{prop:mp-psh-decreasing-stability}
Let $I$ be a directed set and let $(\phi_j)_{j\in I}$ be a
decreasing net in $\PSH(\alpha)$.  If
\[
                         \phi:=\inf_j\phi_j
\]
is not identically $-\infty$, then $\phi\in\PSH(\alpha)$.
\end{proposition}

\begin{proof}
Put $K=X^{\bethan}$ and
$\mathcal P=\PL(K)\cap\PSH(\alpha)$.  For every $j$, choose a
decreasing net
\[
                    h_{j,\lambda}\downarrow\phi_j,
                    \qquad h_{j,\lambda}\in\mathcal P.
\]
Let $\mathcal S$ be the set of all $h\in\mathcal P$ for which
\[
                  h\geq\phi_j+\epsilon
\tag{60.16a}\label{eq:mp-strict-model-dominant}
\]
for some $j\in I$ and some $\epsilon\in\mathbf Q_{>0}$.  This set is
nonempty because rational translation preserves the model psh class.

We claim that $\mathcal S$ is downward directed for pointwise order.
Suppose that $h_i\geq\phi_{j_i}+\epsilon_i$, $i=1,2$, as in
\eqref{eq:mp-strict-model-dominant}.  Choose
$j\in I$ dominating $j_1$ and $j_2$, and choose a positive rational number
\[
                  4\delta<\min\{\epsilon_1,\epsilon_2\}.
\]
The continuous function $g=\min\{h_1,h_2\}$ then satisfies
$\phi_j+4\delta<g$ on the compact space $K$.  The increasing open
sets
\[
                  \{h_{j,\lambda}+2\delta<g\}
\]
cover $K$.  Compactness, followed by directedness of the approximating
net, gives one index $\lambda$ for which this set is all of $K$.
Hence
\[
                  h:=h_{j,\lambda}+\delta\in\mathcal S,
                  \qquad h\leq h_1,\ h_2.
\]
The minimum $g$ is only a continuous comparison function; no closure of
the psh class under minima is being used.

Order $\mathcal S$ by reverse pointwise order.  The identity net is then
a decreasing net in $\mathcal P$.  Every member of $\mathcal S$
dominates $\phi$, so its infimum is at least $\phi$.  Conversely, fix
$x\in K$ and $b>\phi(x)$.  Choose $j$ with $\phi_j(x)<b$, then
choose $\epsilon\in\mathbf Q_{>0}$ and $\lambda$ so that
$h_{j,\lambda}(x)+\epsilon<b$.  Then
$h_{j,\lambda}+\epsilon\in\mathcal S$ is smaller than $b$ at
$x$.  This proves
\[
                         \inf_{h\in\mathcal S}h=\phi
\]
pointwise, including at points where $\phi=-\infty$.  Since the limit is
not identically $-\infty$, \Cref{def:mp-general-psh} gives
$\phi\in\PSH(\alpha)$.
\end{proof}

The class $\PSH(\alpha)$ is stable under arbitrary real translations.
Indeed, approximate a real constant from above by decreasing rational
constants and use the product directed set together with
\Cref{prop:mp-model-psh-stability}.  This observation does not turn an
irrational translate of a rational PL function into a rational PL function;
it applies only after passage to the monotone completion.

The same construction gives decreasing approximation by
$\mathcal H(\alpha)$ for every $\phi\in\PSH(\alpha)$.
Indeed, start with the rational model net in
\Cref{def:mp-general-psh}, approximate each of its terms by
\Cref{prop:mp-real-model-approximation}, and repeat the proof with
$\mathcal H(\alpha)$ in place of the auxiliary class $\mathcal P$.
The directedness argument uses rational translation and the specified
approximating sequences on compact $X^{\bethan}$.

\begin{theorem}\label{thm:mp-psh-divisorial}
Let \(\phi\in\PSH(\alpha)\).
\begin{enumerate}
\item \(\phi(v)>-\infty\) for every \(v\in X^{\divis}\).
\item If \(\psi:X^{\bethan}\to\mathbf R\cup\{-\infty\}\) is upper
      semicontinuous, then
      \[
        \phi\leq\psi\text{ on }X^{\bethan}
        \quad\Longleftrightarrow\quad
        \phi\leq\psi\text{ on }X^{\divis}.
      \tag{60.17}\label{eq:mp-divisorial-comparison}
      \]
\end{enumerate}
\par
\end{theorem}

\begin{proof}[Imported finite-dimensional estimate and deduction]
In dimension zero, $X$ and $X^{\bethan}$ each consist of one point,
so both assertions follow from the definition.  Assume $n\geq1$.
For the first assertion, subtract a constant so that $\sup\phi=0$, and
choose \(\phi_\lambda\downarrow\phi\) as in
\eqref{eq:mp-psh-decreasing-net}.  Replacing each approximant by
\(\phi_\lambda-\sup\phi_\lambda\) does not affect its model positivity, and
the normalized net still converges pointwise to \(\phi\).  Here
\(\sup\phi_\lambda\downarrow\sup\phi=0\): otherwise compactness applied to
the decreasing closed superlevel sets would contradict pointwise
convergence.  Each supremum is a rational model value, so the normalized
approximant remains in \(\PL\).  Fix
\(v=v_{E_0}\in X^{\divis}\) and choose an snc model with central fiber
\(\sum_{j=0}^N b_jE_j\) containing that vertex.  The imported
finite-dimensional estimate of \cite[Theorem~4.3.3, arXiv version~2]{MP24} states that for
every normalized model \(\alpha\)-psh function \(u\),
\[
                  |u(v_{E_j})|\leq C(\cX,\alpha,\gamma)
                  \quad(0\leq j\leq N),
\]
where \(\gamma\) is any fixed K\"ahler class on \(\cX\).  Its proof
intersects the pullback of \(E_j\) with the relatively nef class representing
\(u\) and \(\gamma^{n-1}\); the relations
\(E_j\cdot\cX_0\cdot\gamma^{n-1}=0\), positivity of the off-diagonal
intersections, and connectedness of the central fiber give the uniform
bound.  Passing to the limit at \(v_{E_0}\) proves
\(\phi(v)>-\infty\).

For the second assertion, first suppose that \(\psi\in\PL\).
\cite[Proposition~4.3.6, arXiv version~2]{MP24} gives a finite set
\(\Sigma(\psi)\subset X^{\divis}\), depending only on \(\psi\), such that
\[
             \sup_{X^{\bethan}}(u-\psi)
             =\max_{\Sigma(\psi)}(u-\psi)
             \quad\text{for every }u\in\PSH(\alpha).
\]
Hence the divisorial inequality implies \(\phi\leq\psi\) everywhere.
If $\psi$ is continuous, choose $h_j\in\PL$ converging uniformly
to $\psi$ and rational numbers $q_j>\|h_j-\psi\|_\infty$ with
$q_j\to0$.  Then $h_j+q_j$ remains in $\PL$, and
$\phi\leq\psi<h_j+q_j$ on the divisorial locus.  The PL case gives
$\phi\leq h_j+q_j$ everywhere; taking the uniform limit proves
$\phi\leq\psi$ on $X^{\bethan}$.

Finally, every extended-real upper semicontinuous function on a compact
Hausdorff space is the pointwise decreasing limit of a net of continuous
real functions \(\psi_\lambda\).  If \(\phi\leq\psi\) on
\(X^{\divis}\), then \(\phi\leq\psi_\lambda\) there for every \(\lambda\),
so the continuous case gives \(\phi\leq\psi_\lambda\) everywhere.  Taking
the decreasing limit proves \(\phi\leq\psi\).  The reverse implication is
immediate.  This is \cite[Theorem~4.3.8, arXiv version~2]{MP24}.
\end{proof}

\begin{warning}{Scope of divisorial determination}
In \eqref{eq:mp-divisorial-comparison}, the function \(\phi\) on the left is
\(\alpha\)-psh.  Density of \(X^{\divis}\) alone does not determine an arbitrary
upper semicontinuous function: upper semicontinuity gives inequalities in
the wrong direction.  We use the theorem only with its stated hypothesis.
\end{warning}

For arbitrary psh inputs the mixed energy pairing is defined by
\[
 (\alpha_0,\phi_0)\cdots(\alpha_n,\phi_n)
 :=\inf_{\substack{\psi_i\in\mathcal H(\alpha_i)\\
                    \psi_i\geq\phi_i}}
       (\alpha_0,\psi_0)\cdots(\alpha_n,\psi_n)
 \in\mathbf R\cup\{-\infty\}.
\tag{60.18}\label{eq:mp-extended-pairing}
\]
The use of dominants, rather than an arbitrarily chosen approximating net,
makes the definition canonical.  It is increasing in each potential,
upper semicontinuous, and continuous along decreasing nets
\cite[Lemma~4.4.1 and the discussion following it, arXiv version~2]{MP24}.

The polynomial energy in \Cref{def:mp-energy-j} is defined on
$\PL_{\mathbf R}$ using model representatives.  The extended pairing
now defines energy at arbitrary psh potentials, which need not be model
functions.

\begin{definition}[Energy and \(J\) at arbitrary psh potentials]
\label{def:mp-energy-extended}
For \(\phi\in\PSH(\alpha)\), set
\begin{align*}
 E_\alpha(\phi)
   &:=\frac1{(n+1)V_\alpha}(\alpha,\phi)^{n+1}
     \in\mathbf R\cup\{-\infty\},\\
 J_\alpha(\phi)
   &:=V_\alpha^{-1}(\alpha,\phi)(\alpha,0)^n-E_\alpha(\phi),
\end{align*}
the pairings being the extended ones of
\eqref{eq:mp-extended-pairing}.
\end{definition}

On the real model $\alpha$-psh class this agrees with
\Cref{def:mp-energy-j}.  More generally, let every $\phi_i$ in
\eqref{eq:mp-extended-pairing} be a real model $\alpha_i$-psh function.
The mixed testing formula \eqref{eq:mp-mixed-testing} and nefness show
that every tuple of dominants has pairing at least the model pairing
of the $\phi_i$.  Conversely, \Cref{prop:mp-real-model-approximation}
gives dominants with coefficients converging on fixed models.  On one
common domination of these finitely many models, their intersection
polynomials converge to the model pairing.  The infimum therefore
equals \eqref{eq:mp-mixed-pairing}.

The first term defining $J_\alpha$ is finite.  The zero model potential
has Monge--Amp\`ere measure $\delta_{v_{\triv}}$, so model testing
and decreasing approximation give
\[
 V_\alpha^{-1}(\alpha,\phi)(\alpha,0)^n
       =\phi(v_{\triv})\in\mathbf R.
\]
Finiteness follows from \Cref{thm:mp-psh-divisorial}.  Hence
$J_\alpha(\phi)\in\mathbf R\cup\{+\infty\}$ is well defined even
when $E_\alpha(\phi)=-\infty$.  For all subsequent psh inputs,
$E_\alpha$ and $J_\alpha$ refer to
\Cref{def:mp-energy-extended}.

\begin{definition}[Finite energy]
\label{def:mp-e1}
The finite-energy class is
\[
 \mathcal E^1(\alpha)
 :=\{\phi\in\PSH(\alpha)\mid E_\alpha(\phi)>-\infty\}.
\tag{60.19}\label{eq:mp-e1}
\]
Potentials are considered modulo constants when a normalized space is
needed.
\end{definition}

Mixed pairings of finite-energy inputs are finite by
\cite[Proposition~4.4.4, arXiv version~2]{MP24}.  This permits subtraction
of such pairings.  In particular, for
$\beta\in H^{1,1}_{\mathrm{BC}}(X,\mathbf R)$ and
$\phi\in\mathcal E^1(\alpha)$, choose $C>0$ such that
$C\alpha+\beta$ is K\"ahler and set
\[
 (\beta,0)(\alpha,\phi)^n
 :=(C\alpha+\beta,0)(\alpha,\phi)^n
      -(C\alpha,0)(\alpha,\phi)^n.
\]
The twisted energy is
\[
 E^\beta_\alpha(\phi)
 :=\frac1{V_\alpha}(\beta,0)(\alpha,\phi)^n.
\tag{60.20}\label{eq:mp-twisted-energy}
\]

The two terms are finite.  To check independence of $C$, choose a
decreasing model net $\phi_\lambda\downarrow\phi$.  On each model the
difference is $(\beta,0)(\alpha,\phi_\lambda)^n$, by multilinearity.
Decreasing continuity of the two finite pairings gives the same limit
for every admissible $C$.  This is the extension by differences in
\cite[Definition~4.4.6, arXiv version~2]{MP24}.

The test-function pairing $(0,\psi)(\alpha,\phi)^n$ requires one more
extension when $\psi\in\PL$: an arbitrary model function need not be
psh in any multiple of the pulled-back background on its given model.
Choose a smooth snc model representing $\psi=\phi_D$ and a rational
vertical divisor $A$ such that $\alpha_{\cX}+[A]$ is relatively K\"ahler,
as in \Cref{prop:mp-real-model-approximation}.  For a sufficiently large
positive integer $N$, both
\[
 u=N\phi_A+\psi,\qquad w=N\phi_A
\]
belong to $\mathcal H(N\alpha)$.  Indeed, on each of the finitely many
central components a sufficiently large multiple of the K\"ahler class
$\alpha_{\cX}+[A]$ absorbs the fixed class $[D]$; the other fibres carry
$N\alpha$.  Put $\gamma=N\alpha$ and define
\[
 (0,\psi)(\alpha,\phi)^n
 :=(\gamma,u)(\alpha,\phi)^n
      -(\gamma,w)(\alpha,\phi)^n.
\]
Both terms are finite because $u,w$ are model psh functions.  For any
decreasing model net $\phi_\lambda\downarrow\phi$, model testing and
decreasing continuity give
\[
 V_\alpha^{-1}(0,\psi)(\alpha,\phi)^n
   =\lim_\lambda\int_{X^{\bethan}}\psi\,
                         \MA_\alpha(\phi_\lambda).
\]
For a fixed net the right side depends only on $\psi$, so the definition
does not depend on $A$, $N$, or the chosen representation of $\psi$.
For a fixed representation, the left side is expressed by the canonical
extended pairings, so the limit does not depend on the decreasing net.
Finiteness of both terms is essential before taking their difference.

\begin{proposition}
\label{prop:mp-e1-ma-existence}
For every $\phi\in\mathcal E^1(\alpha)$ there is a unique Radon
probability measure $\MA_\alpha(\phi)$ such that
\[
 \int_{X^{\bethan}}\psi\,\MA_\alpha(\phi)
   =V_\alpha^{-1}(0,\psi)(\alpha,\phi)^n,
 \qquad \psi\in\PL(X^{\bethan}).
\tag{60.21}\label{eq:mp-e1-ma}
\]
If $\phi_\lambda\downarrow\phi$ is any decreasing model psh net, then
$\MA_\alpha(\phi_\lambda)$ converges weakly to $\MA_\alpha(\phi)$.
\end{proposition}

\begin{proof}
Denote the right side of \eqref{eq:mp-e1-ma} by $L_\phi(\psi)$.
The preceding limit formula expresses it as a limit of integrals
against model probability measures.  Hence $L_\phi$ is rational-linear,
is nonnegative on nonnegative model functions, and satisfies
\[
       L_\phi(1)=1,\qquad
       |L_\phi(\psi)|\leq\|\psi\|_\infty.
\]
Uniform density of $\PL$ therefore extends it uniquely to a continuous
real-linear functional on $C^0(X^{\bethan})$.  This extension is positive:
for a nonnegative continuous function, choose uniformly approximating
PL functions and shift them upward by rational bounds for their errors.
Riesz representation on compact $X^{\bethan}$ gives a Radon
probability measure, and density of $\PL$ makes it unique.  This gives
the construction asserted in
\cite[Remark~4.4.8, arXiv version~2]{MP24}.

Finally, for $f\in C^0(X^{\bethan})$ and $\psi\in\PL$, the probability
normalizations give
\[
 \left|\int f\,\MA_\alpha(\phi_\lambda)
               -\int f\,\MA_\alpha(\phi)\right|
 \leq 2\|f-\psi\|_\infty
       +\left|\int\psi\,\MA_\alpha(\phi_\lambda)
               -L_\phi(\psi)\right|.
\]
First choose $\psi$ with arbitrarily small uniform error, and then let
$\lambda$ tend along its directed set.  The limit formula makes the
last term tend to zero, proving weak convergence.
\end{proof}

\section{Radial limits of complex psh rays}

Fix a K\"ahler form $\omega\in\alpha$.  A psh ray is a family
$U=(U_s)_{s>0}$ in $\PSH(X,\omega)$ for which the associated
$S^1$-invariant function
\[
 U(x,\tau):=U_{-\log|\tau|^2}(x),
 \qquad s=-\log|\tau|^2,
\]
is jointly $p_1^*\omega$-psh on $X\times\mathbb D^*$.
This is the subgeodesic condition of \Cref{def:rays-subgeodesic}, in
the radial variable used here; compare
\cite[Section~5.1.1, arXiv version~2]{MP24}.
The slice condition alone is insufficient: the spatially constant
family $U_s=\sin s$ has psh slices, but its radial function is not
subharmonic because $s\mapsto\sin s$ is not convex.

We say that the ray has linear growth if there are $a\geq0$, $b\in\mathbf R$
and $s_0>0$ such that
\[
                      \sup_XU_s\leq as+b
                      \qquad(s\geq s_0).
\tag{60.22}\label{eq:mp-ray-linear-growth}
\]
All constructions concern the germ at $\tau=0$, so we may restrict to
the smaller disc determined by $s_0$.  The function
\[
              \widehat U=U+a\log|\tau|^2
\]
is $p_1^*\omega$-psh on the punctured product and bounded above near
the central fibre.  Its upper semicontinuous extension is therefore
$p_1^*\omega$-psh across that fibre.

For $v\in X^{\divis}$, represent its Gauss extension on a smooth
model $\mu:\cX\to X\times\mathbf P^1$ by
$b_E^{-1}\ord_E$, where $b_E=\ord_E(\tau)$.  Define
\begin{align}
 \nu_v(U)&:=b_E^{-1}\nu_E(\widehat U\circ\mu)-a,
 \tag{60.22a}\label{eq:mp-normalized-ray-lelong}\\
 U^{\NA}(v)&:=-\nu_v(U).
 \tag{60.23}\label{eq:mp-ray-na-value}
\end{align}
where $\nu_E$ is the usual nonnegative generic Lelong number, normalized
by $\nu_E(\log|f|^2)=\ord_E(f)$.  The first term is finite.
Replacing $a$ by another admissible growth bound changes that term by
exactly the change in $a$, so the difference is independent of the bound.
Generic divisorial Lelong numbers are invariant under common domination,
which proves independence of the representing model as well.  The
normalized number $\nu_v(U)$ can be negative.
For the affine ray $U_s=cs$ with $c\in\mathbf R$, one has
$\widehat U=(a-c)\log|\tau|^2$ and
$\nu_E(\widehat U\circ\mu)=(a-c)b_E$.  Thus
$U^{\NA}\equiv c$, including when $c>0$.  More singular rays have
larger normalized Lelong numbers and smaller non-Archimedean potentials.

The squared-norm convention of \Cref{sec:ch51-squared-norm} fixes both
the radial variable and the Lelong normalization in
\eqref{eq:mp-normalized-ray-lelong}.  The factor $b_E^{-1}$ fixes the
Gauss normalization $v(\tau)=1$.  A source using
$s'=-\log|\tau|$ has $s=2s'$; its norm weights and generic Lelong
numbers must be converted together with the radial variable.  The
affine-ray calculation above checks the resulting normalization.

\begin{theorem}\label{thm:mp-ray-to-na}
For every psh ray of linear growth, the function
\eqref{eq:mp-ray-na-value} on $X^{\divis}$ extends uniquely to an element
of $\PSH(\alpha)$.
\end{theorem}

\begin{proof}[Imported regularization input and deduction]
We use the analytic regularization input in the proof of
\cite[Theorem~5.1.4, arXiv version~2]{MP24}.  We spell out the
truncation and limiting steps, since a single fixed logarithmic cutoff
does not retain arbitrarily large divisorial Lelong numbers.

After the growth correction above, subtract a constant and write
$V=\widehat U-b\leq0$ on a smaller product disc.  For each positive
integer $k$, put
\[
       V^{(k)}=\max\{V,k\log|\tau|^2\},
       \qquad \mathfrak a_{k,m}=\mathcal I(mV^{(k)}).
\]
The truncation is jointly $p_1^*\omega$-psh and locally bounded off the
central fibre.  Our multiplier ideals use the weight
$\exp(-mV^{(k)})$.  Since $V^{(k)}\geq k\log|\tau|^2$,
\[
        (\tau^{mk})\subseteq\mathfrak a_{k,m}\subseteq\mathcal O.
\]
Indeed, $|\tau|^{2mk}\exp(-mV^{(k)})\leq1$ is locally integrable.
These invariant coherent ideals determine flag ideals by their finite
expansions modulo $(\tau^{mk})$.  The flag-ideal dictionary therefore
gives rational model functions
\[
        h_{k,m}(v)=-m^{-1}\sigma(v)(\mathfrak a_{k,m}),
        \qquad -k\leq h_{k,m}\leq0,
\]
where $\sigma(v)$ denotes the normalized Gauss extension.

The imported analytic assertion is the following: for each fixed $k$
and $\lambda>1$, and all sufficiently large $m$, regularization on a
smaller disc gives a jointly $p_1^*(\lambda\omega)$-psh potential with
analytic singularity type
$m^{-1}\log|\mathfrak a_{k,m}|^2$.
On a log resolution with
$\mathfrak a_{k,m}\mathcal O_{\mathcal Y}=\mathcal O_{\mathcal Y}(-F_{k,m})$,
the bounded residual potential gives relative nefness of
$\lambda\alpha_{\mathcal Y}-m^{-1}[F_{k,m}]$; this is the
bounded-potential argument of
\cite[Lemma~5.1.3, arXiv version~2]{MP24}.
Thus $h_{k,m}$ is model $(\lambda\alpha)$-psh.  The regularization
input, including its curvature control, remains imported here.

Subadditivity gives
$\mathfrak a_{k,2m}\subseteq\mathfrak a_{k,m}^{2}$, and hence
$h_{k,2m}\leq h_{k,m}$.  For the dyadic sequence this produces a limit
$h_k$, bounded between $-k$ and $0$ on all of $X^{\bethan}$.
For every $\lambda>1$ a tail lies in
$\PL\cap\PSH(\lambda\alpha)$, so $h_k\in\PSH(\lambda\alpha)$.
The multiplier-ideal Lelong formula
\eqref{eq:complex-lelong-from-ideals}, applied on the product, gives
\[
 h_k(v)
   =-b_E^{-1}\nu_E(V^{(k)}\circ\mu)
   =\max\{-b_E^{-1}\nu_E(V\circ\mu),-k\}
       \qquad(v\in X^{\divis}).
\]
In particular, the cutoff changes the divisorial limit whenever the
normalized Lelong number exceeds $k$.

To return to the original class, take
$\lambda_j=1+1/j\downarrow1$.  Pointwise division of the model approximants
by $\lambda_j$ rescales their background class to $\alpha$, so
$h_k/\lambda_j\in\PSH(\alpha)$.  Since $h_k\leq0$, these functions
decrease to $h_k$.  \Cref{prop:mp-psh-decreasing-stability} gives
$h_k\in\PSH(\alpha)$.  This division is ordinary multiplication of
potentials with simultaneous rescaling of the background, not the
valuation-scaling action \eqref{eq:mp-scaling-action}.

Finally, $V^{(k+1)}\leq V^{(k)}$ gives inclusions of multiplier ideals,
so $h_{k+1}\leq h_k$ on the whole spectrum.  Their limit $h$ is not
identically $-\infty$, since its value at $v_{\triv}$ is the negative
of the finite generic Lelong number of $V$ along $X\times\{0\}$.
Decreasing stability yields $h\in\PSH(\alpha)$, and the displayed
divisorial formula shows that $h+a$ has the values
\eqref{eq:mp-ray-na-value}.  Real translation preserves $\PSH(\alpha)$,
so this gives the required extension.  Uniqueness follows from
\Cref{thm:mp-psh-divisorial}.
\end{proof}

Conversely, a finite-energy potential determines a ray only after a complex
initial potential is chosen.  We first specify the comparison property.

\begin{definition}\label{def:mp-maximal-ray}
A finite-energy geodesic ray $U$ is maximal if, for every finite-energy
geodesic ray $V$ of linear growth,
\[
 V_0\leq U_0,\qquad V^{\NA}\leq U^{\NA}
 \quad\Longrightarrow\quad
 V_s\leq U_s\quad(s\geq0).
\]
Here the initial values are the finite-energy boundary values.
This is the geodesic comparison in
\cite[Definition~5.1.6, arXiv version~2]{MP24}.
\end{definition}

\begin{theorem}\label{thm:mp-maximal-ray}
For $\phi\in\mathcal E^1(\alpha)$ and
$u\in\mathcal E^1(X,\omega)$, there is a unique maximal finite-energy
geodesic ray $U$ with $U_0=u$ and $U^{\NA}=\phi$.
\end{theorem}

\begin{proof}[Imported existence and direct uniqueness]
Existence, including the identification of the non-Archimedean data of
the limiting ray, is \cite[Theorem~5.1.7, arXiv version~2]{MP24}.
Its model construction uses a dominating model on which the represented
class is relatively K\"ahler; see
\cite[Lemma~5.1.8]{MP24} in the same version.  The general finite-energy
construction and the passage from model data to the prescribed
$\phi$ are included in this imported assertion.

Uniqueness follows directly from \Cref{def:mp-maximal-ray}.  If $U$ and
$V$ have the same initial value and the same non-Archimedean data, the
maximality of $U$ gives $V\leq U$, and that of $V$ gives $U\leq V$.
\end{proof}

\begin{remark}\label{rem:mp-maximal-limit-inputs}
The energy calculation in the model approximation must be distinguished
from the identification of the limiting non-Archimedean potential.
Suppose that decreasing model data $\phi_m\downarrow\phi$ and smooth
initial values $u_m\downarrow u$ have been chosen in the construction,
with associated maximal rays $U^m$.  The model identity
\cite[Remark~5.1.9, arXiv version~2]{MP24} and monotone energy continuity
give, for each fixed $s$,
\[
 E_\omega(U^m_s)
   =E_\omega(u_m)+sE_\alpha(\phi_m)
   \downarrow E_\omega(u)+sE_\alpha(\phi).
\]
The finite lower bound controls the complex energy of the decreasing
limit $U_s$.  Order gives $U^{\NA}\leq\phi$.  These observations alone
do not identify $E_\alpha(U^{\NA})$.

In the projective polarized argument,
\cite[Theorem~6.4, arXiv version~3]{BBJ21} supplies the additional bound
\[
 E_\alpha(U^{\NA})
 \geq\lim_{s\to\infty}s^{-1}E_\omega(U_s),
\]
and \cite[Lemma~6.3]{BBJ21} in the same version supplies strict
monotonicity on ordered finite-energy potentials.  Together with
$U^{\NA}\leq\phi$, these give
\[
 E_\alpha(\phi)\geq E_\alpha(U^{\NA})
 \geq\lim_{s\to\infty}s^{-1}E_\omega(U_s)=E_\alpha(\phi).
\]
Equality of the energies then gives $U^{\NA}=\phi$.
This is the identification in the proof of its Theorem~6.6, and it
uses more than monotone continuity of the model energy.
For the general K\"ahler assertion above, the identification remains
part of the imported MP existence theorem; the projective theorem is
not being applied to an arbitrary transcendental class.
Likewise, replacing the model nets of
\eqref{eq:mp-psh-decreasing-net} by a decreasing sequence is an additional
regularization step in the construction, not a consequence of the
definition of $\PSH(\alpha)$.
\end{remark}

\begin{proposition}\label{prop:mp-maximal-order}
Write $U(u,\phi)$ for the ray in \Cref{thm:mp-maximal-ray}.
If $u\leq v$ and $\phi\leq\psi$, then
$U(u,\phi)_s\leq U(v,\psi)_s$ for every $s\geq0$.
For $b,c\in\mathbf R$ one has
\[
 U(u+b,\phi+c)_s=U(u,\phi)_s+b+cs.
\]
In particular, if $u-C\leq v\leq u+C$ for $C\geq0$, then
\[
 U(u,\phi)_s-C\leq U(v,\phi)_s\leq U(u,\phi)_s+C
 \qquad(s\geq0).
\]
\end{proposition}

\begin{proof}
The first assertion is the comparison property of $U(v,\psi)$.
Adding $b+cs$ preserves the geodesic equation and changes the initial
value by $b$.  By \eqref{eq:mp-normalized-ray-lelong}, it changes the
non-Archimedean potential by $c$; the constant $b$ has no divisorial
Lelong number.  To check maximality, subtract $b+cs$ from any competing
ray and apply the maximality of $U(u,\phi)$.  Uniqueness then proves
the translation formula.  Applying order preservation to
$u-C\leq v\leq u+C$ and using that formula with $c=0$ proves the last
assertion.
\end{proof}

\begin{example}\label{ex:mp-ray-limit-loss}
The map $U\mapsto U^{\NA}$ and the energy slope need not commute
with decreasing limits of psh rays.  On any compact K\"ahler manifold
consider the spatially constant rays
\[
 U^j_s=-\min(s,j),\qquad U_s=-s,\qquad j\geq1.
\]
Their complexifications are
$\max\{\log|\tau|^2,-j\}$ and $\log|\tau|^2$, respectively.
Thus they are joint psh rays with initial value zero and finite-energy
slices, and $U^j_s\downarrow U_s$ for every fixed $s$.
Each $U^j$ is bounded near the central fibre, whereas $U$ has normalized
Lelong number one there.  Consequently
\[
 (U^j)^{\NA}=0,\qquad U^{\NA}=-1.
\]
With $E_\omega(0)=0$, translation of the energy gives
$E_\omega(U^j_s)=-\min(s,j)$ and $E_\omega(U_s)=-s$.  In particular,
\[
 \lim_{j\to\infty}\lim_{s\to\infty}
       \frac{E_\omega(U^j_s)}s=0,
 \qquad
 \lim_{s\to\infty}\lim_{j\to\infty}
       \frac{E_\omega(U^j_s)}s=-1.
\]
The approximating rays have a corner at $s=j$ and are not geodesic
rays.  The example concerns general psh rays and does not contradict
\Cref{thm:mp-maximal-ray}.  It shows why a decreasing limit alone cannot
supply the energy identification in that theorem.
\end{example}

\section{Comparison with the test-curve construction}

\Cref{chap:transcendental-metrics} follows the Darvas--Xia--Zhang
test-curve language.  On a compact K\"ahler manifold the following
finite-energy bridge is available.  In the connected smooth projective
ample N\'eron--Severi setting it is complemented by the full cone comparison
of \Cref{thm:pseff-projective-full-mp-bj}.

\begin{theorem}\label{thm:mp-dxz-common-locus}
Let \(\beta\) be a K\"ahler class.  On the finite-energy class of
\(\mathcal I\)-maximal test curves of \cite{DXZ23}, the Ross--Witt
Nystr\"om transform gives maximal geodesic rays, and taking radial
non-Archimedean data gives a bijection with
\(\mathcal E^1(\beta)\).  Under this correspondence, the energy and mixed
energy slopes agree.
\end{theorem}

\begin{proof}[Imported common-locus correspondence]
The test-curve-to-ray statement is the finite-energy Ross--Witt Nystr\"om
correspondence of \cite[Theorem~3.17]{DXZ23}, used in
\cite[Section~5.2]{MP24}.
\Cref{thm:mp-maximal-ray} gives the inverse after fixing the same initial
potential.  Equality of the energies is a special case of the slope formula
\cite[Theorem~5.3.4]{MP24}.  No claim is made here for arbitrary big
classes, nonmaximal rays, or test curves outside the finite-energy locus.
\end{proof}

\begin{warning}{Scope of the full comparison}
The DXZ construction starts from singularity types and test curves on the
complex manifold; the MP construction starts from coherent ideals and
analytic models.  The comparison theorem above identifies the objects only
after the hypotheses K\"ahler, \(\mathcal I\)-maximal, and finite energy are
imposed, and it additionally identifies radial energies.  There is one
full-cone extension: if \(X\) is connected and smooth projective and
\(\alpha\in\operatorname{Amp}(X)\subset N^1(X)_{\mathbf R}\), then
\Cref{thm:pseff-projective-full-mp-bj} identifies the unrestricted MP and
BJ cones.  When \(\alpha=c_1(L)\) for an ample \(\mathbf Q\)-line bundle,
composition with \cite[Theorem~3.14]{DXZ23} identifies the unrestricted
DXZ and MP cones on all semivaluations, with no finite-energy hypothesis.

This full comparison does not extend here to a nonprojective compact
K\"ahler manifold, to a transcendental Bott--Chern class outside the real
N\'eron--Severi space, or to a non-nef class.  It is an order isomorphism of
psh cones; compatibility of energy, mixed pairings, and Monge--Amp\`ere
operators requires separate arguments and is not inferred from cone
equality.
\end{warning}

\section{Slope formulas and stability functionals}

Let \(\eta_i\) be smooth closed real \((1,1)\)-forms.  If
\(\phi_i\in\PL\), choose smooth rays \(U_i\) compatible with the
corresponding analytic models.  If
\(\phi_i\in\mathcal E^1(\alpha)\), choose the maximal rays of
\Cref{thm:mp-maximal-ray}.

\begin{theorem}\label{thm:mp-mixed-slope}
Suppose
\[
 \phi_0,\ldots,\phi_k\in\PL,
 \qquad
 \phi_{k+1},\ldots,\phi_n\in\mathcal E^1(\alpha).
\]
With compatible rays as above,
\[
 \lim_{s\to+\infty}\frac1s
   (\eta_0,U_{0,s})\cdots(\eta_n,U_{n,s})
   =([\eta_0],\phi_0)\cdots([\eta_n],\phi_n).
\tag{60.24}\label{eq:mp-mixed-slope}
\]
\end{theorem}

\begin{proof}[Imported quantitative mixed-energy core]
The imported result is \cite[Theorem~5.3.4]{MP24}, whose hypotheses and
normalization are those in the statement.  When every input is PL, it is
the cohomological intersection formula for
smooth compatible rays.  Approximate each finite-energy input from above by
\(\mathcal H(\alpha)\)-potentials.  Monotone continuity of the extended
pairing gives convergence on the non-Archimedean side.  On the complex side,
the quantitative mixed-energy estimate controls the error by a fractional
power of the \(d_1\)-distance.  Maximality turns this distance into an energy
difference whose asymptotic slope tends to zero.  The quantitative estimate
has exponent \(2^{-n}\); after division by \(s\), its right side is a
fractional power of an energy-slope error and therefore tends to zero.
Iterating over the finite-energy inputs proves the formula.
\end{proof}

In particular, for the maximal ray \(U\) associated to \(\phi\),
\[
 \lim_{s\to\infty}\frac{E_\omega(U_s)}s=E_\alpha(\phi),
 \quad
 \lim_{s\to\infty}\frac{E^\eta_\omega(U_s)}s
       =E^{[\eta]}_\alpha(\phi),
 \quad
 \lim_{s\to\infty}\frac{J_\omega(U_s)}s=J_\alpha(\phi).
\tag{60.25}\label{eq:mp-basic-slopes}
\]

Define the reduced entropy and non-Archimedean Mabuchi functional by
\[
 \Ent_\alpha^{\NA}(\phi)
   :=\int_{X^{\bethan}}A_X^{\NA}\,\MA_\alpha(\phi),
\qquad
 M_\alpha^{\NA}
   :=\bar S_\alpha E_\alpha+E_\alpha^{-c_1(X)}
     +\Ent_\alpha^{\NA},
\tag{60.26}\label{eq:mp-na-mabuchi}
\]
where
\[
                   \bar S_\alpha
                  :=n\,\frac{c_1(X)\cdot\alpha^{n-1}}{\alpha^n}.
\tag{60.27}\label{eq:mp-average-scalar}
\]
Our \(A_X^{\NA}=A^{\mathrm{tot}}-1\) is the relative discrepancy convention
of \Cref{chap:mp-dual-complexes}.  Replacing it by \(A^{\mathrm{tot}}\) adds
\(1\), since the Monge--Amp\`ere measure has mass one.  Thus, if one
formally defines a total-discrepancy functional
\(M_{\alpha,\mathrm{tot}}^{\NA}\), then
\[
             M_{\alpha,\mathrm{tot}}^{\NA}
             =M_\alpha^{\NA}+1.
\tag{60.27a}\label{eq:mp-total-reduced-mabuchi}
\]
The shift is not harmless for slopes: the complex Mabuchi functional is
invariant under constants, whereas the total-discrepancy expression has
value \(1\) at the zero non-Archimedean potential.  The function constructed
in \cite[(3.2.4)]{MP24} is \(A^{\mathrm{tot}}\), while the reduced
normalization is stated explicitly in \cite[Definition~2.4.1]{MPWN25}.

\begin{warning}{Boundary of the entropy-slope argument}
Proposition~6.2.2 and Corollary~6.2.3 of \cite{MP24} are printed with
\(A^{\mathrm{tot}}\).  Their claimed inequality cannot hold in that
normalization: for \(\phi=0\),
\[
 \MA_\alpha(0)=\delta_{v_{\triv}},\qquad
 M_{\alpha,\mathrm{tot}}^{\NA}(0)=1,
 \qquad
 \lim_{s\to+\infty}\frac{M_\omega(0)}s=0.
\tag{60.28}\label{eq:mp-total-slope-obstruction}
\]
The later paper uses reduced discrepancy, but does not state a separate
corrected entropy-slope theorem.  Replacing the PL approximants in the
printed proof by \(A^{\mathrm{tot}}\circ p_{\cX}-1\) is the expected repair;
without an independent verification of all compatibility and integrability
steps, this book does not promote that repair to a theorem.  Formally, if
the reduced entropy inequality, with \(H_\mu\) denoting complex entropy
relative to the normalized volume measure,
\[
 \Ent_\alpha^{\NA}(\phi)
 \leq\liminf_{s\to+\infty}\frac{H_\mu(U_s)}s
\]
holds for a maximal ray, then adding the first two proved identities in
\eqref{eq:mp-basic-slopes} gives the corresponding Mabuchi inequality.
No later argument in this book uses that formal implication as an
independently proved result.
\end{warning}

\begin{definition}[Uniform stability over finite energy]
\label{def:mp-uniform-e1-stability}
The pair \((X,\alpha)\) is uniformly K-stable over
\(\mathcal E^1\) if some \(\delta>0\) satisfies
\[
            M_\alpha^{\NA}(\phi)\geq\delta J_\alpha(\phi)
       \quad\text{for all }\phi\in\mathcal E^1(\alpha).
\tag{60.29}\label{eq:mp-uniform-e1}
\]
\end{definition}

\begin{theorem}\label{thm:mp-e1-stability-csck}
If a compact K\"ahler pair \((X,\alpha)\) is uniformly K-stable over
\(\mathcal E^1\), then \(\alpha\) contains a unique cscK metric.
\end{theorem}

\begin{proof}[Imported later-source theorem]
This is exactly \cite[Theorem~2.4.3]{MPWN25}.  The normalization check is
part of the import: \cite[Definition~2.4.1]{MPWN25} defines entropy using
\(A_{X\times\mathbf P^1}-1\), and its Definition~2.4.2 is
\eqref{eq:mp-uniform-e1}; the scalar-curvature coefficient is understood
with the independently verified factor recorded in
\eqref{eq:mpwn-s-constant}.  Thus the hypotheses and functional in the
cited theorem coincide with those above, and its conclusion gives the
claimed cscK metric.  This proof treats the later theorem as a deep imported
input; it does not derive it from the unverified normalization repair in the
preceding warning.
\end{proof}

\section*{Source note}

Model psh functions and their closedness are MP24 Section~4.1.  The atomic
Monge--Amp\`ere measure, mixed pairing, and energy are Sections~4.2--4.4.
Their algebraic model-calculus antecedents include
\cite[Proposition~3.6 and Lemma~3.24]{BJ22}; the analytic maximum theorem and
closedness argument are MP inputs and are not inferred from that algebraic
specialization.
The radial construction, maximal rays, comparison with test curves, and
mixed slope theorem are Sections~5.1--5.3.  MP24 Sections~6.2--6.3 contain
the printed entropy-slope argument and cscK implication.  The literal
discrepancy constructed in MP24 Section~3.2 is total, so
\eqref{eq:mp-total-reduced-mabuchi} and
\eqref{eq:mp-total-slope-obstruction} record the exact normalization
boundary.  Reduced entropy and the finite-energy cscK implication are
imported from \cite[Definitions~2.4.1--2.4.2 and
Theorem~2.4.3]{MPWN25}; no independently corrected entropy-slope theorem is
claimed here.
\Cref{chap:mpwn-cy} uses later results of \cite{MPWN25} to
upgrade this theory to a full finite-energy Calabi--Yau correspondence;
\Cref{chap:mpwn-k-stability} explains why that upgrade is also needed before model test
configurations can be used as a complete stability test.
